\FloatBarrier
\section{Conclus\~ao}

Este estudo avaliou a detec\c{c}\~ao e contagem de c\'elulas cervicais a partir de
anota\c{c}\~oes pontuais, tratando a convers\~ao de ponto em caixa como uma decis\~ao
metodol\'ogica expl\'icita. Em resposta \`a RQ1, o sweep explorat\'orio mostrou que o
tamanho da pseudo-caixa afeta F1, AP50 e MAE de forma n\~ao monot\^onica: tamanhos entre
96 e 192 pixels, exceto 160, formam uma regi\~ao de desempenho pr\'oximo, mas com
compromissos distintos de AP50, F1 e MAE de contagem.
O lado de 144 pixels foi selecionado como ponto operacional intermedi\'ario e
levado ao protocolo final.

Em resposta \`a RQ2, na valida\c{c}\~ao cruzada 5-fold, com testes externos por imagem e
limiar selecionado apenas na valida\c{c}\~ao interna de cada fold, o YOLO11s obteve
F1=0,7969$\pm$0,0267, AP50=0,8569$\pm$0,0256 e MAE=5,85$\pm$1,09 c\'elulas por
imagem. Esses valores resumem o desempenho final do
protocolo, estimado sobre todos os testes externos da valida\c{c}\~ao cruzada.

A conclus\~ao principal n\~ao \'e que 144 pixels sejam universais. O resultado
relevante \'e que a geometria derivada de pontos altera a localiza\c{c}\~ao e a
contagem, devendo ser escolhida, reportada e validada. Em termos pr\'aticos, o
detector de classe \'unica pode funcionar como etapa anterior a um classificador
celular ou a uma revis\~ao assistida, desde que falsos positivos, duplicidades e
tend\^encia de supercontagem sejam considerados.

Como limita\c{c}\~oes, as m\'etricas baseadas em IoU descrevem consist\^encia com
pseudo-caixas, n\~ao com fronteiras anat\^omicas; a base n\~ao fornece
identificador de paciente, restringindo a independ\^encia ao n\'ivel de imagem; e
o sweep de tamanho foi explorat\'orio, com uma execu\c{c}\~ao por configura\c{c}\~ao.
Trabalhos futuros devem comparar pseudo-caixas a m\'etodos que aprendem
extens\~ao a partir de pontos \cite{chen2023p2bnet}, repetir a escolha do tamanho
com m\'ultiplas sementes e anotar uma amostra com contornos manuais para calibrar
a interpreta\c{c}\~ao espacial.

O reposit\'orio p\'ublico do projeto\footnote{\url{https://github.com/JeanCarlos899/Pesquisa-Jean}}
re\'une o c\'odigo da pipeline, arquivos de configura\c{c}\~ao, CSVs usados na
gera\c{c}\~ao das tabelas e figuras, al\'em dos resultados da valida\c{c}\~ao cruzada.
Essa organiza\c{c}\~ao permite auditar os valores reportados e reproduzir as
principais etapas experimentais descritas no artigo.
